% !TEX root = ../notes_sovereignai.tex
% Chapter 4: A Voice in the House. HA + Wyoming + Whisper + Piper + Ollama.
%%%%%%%%%%%%%%%%%%%%%%%%%%%%%%%%%%%%%%%%%%%%%%%%%%%%%%%%%%%%%%%%%%%%%%

\index{Home Assistant}\index{Wyoming}\index{Whisper}\index{Piper}\index{voice assistant}
\chapter{A Voice in the House}
\printchapterversion{04_voice}
\newthought{Speak, and It Answers.}
\vspace{1.5em}

% ==============================================================================================
% THE WHY
% ==============================================================================================

\section{The Why}\index{smart home}\index{privacy}

\newthought{The engine reaches every screen you own} and none of the air in the room. A voice
assistant closes that gap, and it is the one piece of household AI most people
already rent in its worst form: an always-listening microphone that ships every
word to a company's servers. The sovereign version hears, thinks, and speaks
entirely on your hardware, and it can act on the house rather than only answer
trivia.

\newthought{The orchestrator is Home Assistant.}\index{Home Assistant!container} It is the open hub the smart-home
world already speaks, it runs as a container on the rig, and its voice pipeline is
built from swappable parts joined by a small protocol. We give it a local ear, a
local voice, and the local model as its brain.

\newthought{By the end of this chapter} you will have:
\begin{itemize}
    \item run Home Assistant with local speech-to-text and text-to-speech;
    \item made the Ollama model its conversation brain, able to act on devices;
    \item spoken to a satellite and heard the rig answer, with nothing leaving the house.
\end{itemize}

% ==============================================================================================
% SPEECH TO MODEL TO SPEECH
% ==============================================================================================

\section{Speech to Model to Speech}\index{speech to text}\index{text to speech}%
\index{STT|see{speech to text}}\index{TTS|see{text to speech}}

\newthought{A voice turn is three steps.} Speech-to-text turns the spoken request
into words; the model decides what it means and what to do; text-to-speech turns
the answer back into sound. The local stack is faster-whisper for the ear, the
Chapter~2 model for the brain, and Piper for the voice, each a small service, each
replaceable.

\marginnote{\newthought{Wyoming}\index{Wyoming!protocol} is Home Assistant's protocol for voice parts: each
service listens on a TCP port and speaks audio and events, and Home Assistant wires
them into a pipeline. Whisper listens on \texttt{10300}, Piper on \texttt{10200}.}

\newthought{Put the ear on the GPU.}\index{faster-whisper} Speech-to-text is the slow step, and the
built-in add-on runs on the CPU; running faster-whisper as its own container with
the GPU passed through brings a request under a second. The same 3090 that serves
chat transcribes speech in the gaps, since the two rarely peak at once in a house.

\marginnote{\newthought{The brain must hold tools.}\index{tool calling}\index{conversation agent} A conversation agent that only
chats can answer the weather but cannot turn off a light. Home Assistant exposes its
devices to the model as tools, so the model must be one that supports tool calling;
the Qwen3 model from Chapter~2 was chosen for exactly this.}

% ==============================================================================================
% WIRE THE PIPELINE
% ==============================================================================================

\section{Wire the Pipeline}\index{voice pipeline}

\newthought{Run the orchestrator and the voice parts} as containers on the rig:
Home Assistant, the ear with the GPU passed through, and the voice.

\begin{verbatim}
$ docker run -d --name homeassistant --network host \
    -v ha-config:/config ghcr.io/home-assistant/home-assistant:stable

$ docker run -d --name whisper -p 10300:10300 --gpus all \
    rhasspy/wyoming-whisper --model small-int8 --language en --device cuda

$ docker run -d --name piper -p 10200:10200 \
    rhasspy/wyoming-piper --voice en_US-lessac-medium
\end{verbatim}

\marginnote{\newthought{Add the parts in the UI.}\index{Home Assistant!Wyoming integration} In Home Assistant, Settings,
Devices and Services, add the Wyoming Protocol integration twice, pointing it
at the rig's address on ports \texttt{10300} and \texttt{10200} to register the ear
and the voice.}

\newthought{Make the model the brain.}\index{Home Assistant!Ollama integration} Add the Ollama integration and point it at
the engine, then choose a tool-calling model:

\begin{verbatim}
Settings > Devices & Services > Add Integration > Ollama
  URL ..... http://localhost:11434
  Model ... qwen3:14b
  [x] let the model control Home Assistant devices
\end{verbatim}

\newthought{Assemble the pipeline.}\index{Home Assistant!voice pipeline} Under Settings, Voice Assistants, build one
assistant from the three local parts: faster-whisper for speech-to-text, the Ollama
agent for conversation, and Piper for text-to-speech. Set it as the preferred
assistant.

\marginnote{\newthought{Give it ears in a room.}\index{voice satellite}\index{Wyoming!satellite} A Home Assistant Voice PE puck, or
any Wyoming satellite, is the microphone and speaker on the shelf; it hands audio to
the pipeline and plays the answer back. One per room you want to talk to.}

\newthought{What you now own} answers out loud. You say ``turn off the kitchen light
and tell me a joke.''
The satellite streams it to the rig, faster-whisper transcribes it, the model turns
off the light through Home Assistant and writes a reply, and Piper says it back, all
on hardware you own, with the router unplugged if you like.

% ==============================================================================================
% BRIDGE
% ==============================================================================================

\newthought{The house can now listen and speak}, and chat, terminal, and voice all
draw on the one engine. What the engine still lacks is your own life: it answers
from what it was trained on, never from what you have written. The next chapter
points your notes at it.

\marginnote{\newthought{feedback}}
